% ============================================================
% Pose-Differentiable Body Twins (RIHB-in-the-loop) — v2
% First-principles companion to rib.tex and q_complement.tex.
% Intended home for the deployable parts: a JSAC2 SI submission
% on Digital Twins for Wireless Networks (Q4 2026 issue).
% Author: R. Wydaeghe
%
% v2 addresses critique_v1.md:
%   - new lemma on dB-per-degree LOS un-shadowing rate gradient (K1, I6)
%   - LOS un-shadowing theorem: rate up AND dose down, both monotone in v
%   - MCS cap saturation plateau corollary (K4)
%   - Pareto bound restated as trace inequality, fixed operator slip (K2, K6)
%   - pose-noise-aware Pareto (K1)
%   - near-field correction for body-to-phone propagation (K3)
%   - gait-action principle restricted to body-mediated component (K5)
%   - convergence under partial compliance (I4)
%   - acceptance-pattern privacy caveat (I3)
%   - 1% opt-in back-of-envelope (I1)
%   - RULA/REBA reference for comfort weights (I2)
%   - notation cleanup (I5)
%   - block diagram, "no new hardware", tightened bandwidth arithmetic (C1, C2, C4)
% ============================================================

\documentclass[11pt,a4paper]{article}

\usepackage[margin=2.4cm]{geometry}
\usepackage[utf8]{inputenc}
\usepackage[T1]{fontenc}
\usepackage{lmodern}
\usepackage{microtype}
\linespread{1.04}

\usepackage{amsmath,amssymb,amsthm,mathtools}
\usepackage{bm}
\usepackage{mathrsfs}
\usepackage{bbm}
\usepackage{cancel}
\usepackage{xfrac}
\usepackage{booktabs}
\allowdisplaybreaks[2]

\usepackage{empheq}
\usepackage[most]{tcolorbox}
\tcbset{highlight math style={%
    enhanced, colframe=black!70, colback=gray!6,
    arc=2pt, boxrule=0.5pt, left=3pt, right=3pt}}

\usepackage{siunitx}
\sisetup{detect-all, range-phrase = {--}, range-units = single,
         per-mode = symbol, inter-unit-product = {\,}}

\usepackage{enumitem}
\usepackage{xcolor}
\usepackage{xspace}
\usepackage{fancyhdr}
\pagestyle{fancy}
\fancyhf{}
\cfoot{\thepage}
\renewcommand{\headrulewidth}{0pt}

\usepackage{titlesec}
\titleformat{\section}{\large\bfseries}{\thesection}{1em}{}
\titleformat{\subsection}{\normalsize\bfseries}{\thesubsection}{1em}{}

\usepackage{tikz}
\usetikzlibrary{positioning,arrows.meta,calc,fit,backgrounds,shapes.geometric}

\usepackage[colorlinks=true, linkcolor={blue!55!black},
            citecolor={green!45!black}, urlcolor={blue!65!black}]{hyperref}
\usepackage[nameinlink,capitalize]{cleveref}

\usepackage{aliascnt}
\theoremstyle{definition}
\newtheorem{theorem}{Theorem}[section]
\newaliascnt{proposition}{theorem}
\newtheorem{proposition}[proposition]{Proposition}
\aliascntresetthe{proposition}
\newaliascnt{lemma}{theorem}
\newtheorem{lemma}[lemma]{Lemma}
\aliascntresetthe{lemma}
\newaliascnt{corollary}{theorem}
\newtheorem{corollary}[corollary]{Corollary}
\aliascntresetthe{corollary}
\newaliascnt{definition}{theorem}
\newtheorem{definition}[definition]{Definition}
\aliascntresetthe{definition}
\newaliascnt{remark}{theorem}
\newtheorem{remark}[remark]{Remark}
\aliascntresetthe{remark}
\crefname{proposition}{Proposition}{Propositions}
\Crefname{proposition}{Proposition}{Propositions}
\crefname{lemma}{Lemma}{Lemmas}\Crefname{lemma}{Lemma}{Lemmas}
\crefname{corollary}{Corollary}{Corollaries}\Crefname{corollary}{Corollary}{Corollaries}
\crefname{definition}{Definition}{Definitions}\Crefname{definition}{Definition}{Definitions}
\crefname{remark}{Remark}{Remarks}\Crefname{remark}{Remark}{Remarks}

% Macros aligned with monograph, q_complement, rib.tex
\newcommand{\khat}{\hat{\bm{k}}}
\newcommand{\nhat}{\hat{\bm{n}}}
\newcommand{\rhat}{\hat{\bm{r}}}
\newcommand{\rr}{\mathbf{r}}
\newcommand{\EE}{\mathbf{E}}
\newcommand{\Sinc}{S_{\mathrm{inc}}}
\newcommand{\Sab}{S_{\mathrm{ab}}}
\newcommand{\Sref}{S_{\mathrm{re}}}
\newcommand{\Sinn}{S_{\mathrm{in}}}
\newcommand{\Smis}{S_{\mathrm{mi}}}
\DeclareMathOperator{\ReLU}{ReLU}
\DeclareMathOperator{\RE}{Re}
\DeclareMathOperator{\IM}{Im}
\DeclareMathOperator{\tr}{tr}
\DeclareMathOperator*{\argmin}{arg\,min}
\DeclareMathOperator*{\argmax}{arg\,max}
\newcommand{\diff}{\mathrm{d}}
\newcommand{\ident}{\mathbbm{1}}
\newcommand{\Reals}{\mathbb{R}}
\newcommand{\Complex}{\mathbb{C}}
\newcommand{\half}{\tfrac{1}{2}}
\newcommand{\Gtil}{\tilde{\mathbf{G}}}
\newcommand{\gtil}{\tilde{\mathbf{g}}}
\newcommand{\Fn}{\mathbf{F}_{\!n}}
\newcommand{\Rn}{\mathbf{R}_{\!n}}
\newcommand{\Qabs}{\mathbf{Q}_{\mathrm{ab}}}
\newcommand{\Qref}{\mathbf{Q}_{\mathrm{re}}}
\newcommand{\Qinn}{\mathbf{Q}_{\mathrm{in}}}
\newcommand{\Qmis}{\mathbf{Q}_{\mathrm{mi}}}

% RIB / RIHB macros
\newcommand{\kin}{\hat{\bm{k}}_{\mathrm{i}}}
\newcommand{\kout}{\hat{\bm{k}}_{\mathrm{o}}}
\newcommand{\Frib}{\mathbf{F}_{\mathrm{b}}}
\newcommand{\sigb}{\sigma_{\mathrm{b}}}
\newcommand{\Qbist}{\mathbf{Q}_{\mathrm{bi}}}
\newcommand{\hh}{\mathbf{h}}
\newcommand{\Jmat}{\mathbf{J}}
\newcommand{\bvec}{\mathbf{b}}
\newcommand{\xvec}{\mathbf{x}}
\newcommand{\thetab}{\bm{\theta}}
\newcommand{\betab}{\bm{\beta}}
\newcommand{\Mmesh}{\mathcal{M}}
\newcommand{\LBS}{\mathrm{LBS}}
\newcommand{\Pabs}{P_{\mathrm{ab}}}
\newcommand{\PLOS}{P_{\mathrm{LOS}}}
\newcommand{\Pbody}{P_{\mathrm{body}}}
\newcommand{\Lrl}{L_{\mathrm{RL}}}
\newcommand{\SINR}{\mathrm{SINR}}

\title{\textbf{Pose-Differentiable Body Twins:\\
First Principles for User-Adaptive\\
mmWave Capacity and Exposure (v2)}}
\author{Robin Wydaeghe\\
 Ghent University \& imec\\
 \texttt{robin.wydaeghe@ugent.be}}
\date{\today}

\makeatletter
\renewcommand{\maketitle}{%
  \thispagestyle{plain}%
  \begin{center}
    \rule{\linewidth}{1pt}\\[0.4em]
    {\LARGE\bfseries \@title}\\[0.5em]
    {\large \@author}\\[0.3em]
    {\small \@date}\\[0.4em]\rule{\linewidth}{1pt}
  \end{center}
  \vspace{0.5em}
}
\makeatother

\begin{document}
\maketitle

\begin{abstract}
A user with a smartphone is, at mmWave frequencies, an electromagnetic
object that the network can both \emph{model} and \emph{negotiate
with}.  We treat the user's pose
$\thetab\in\Reals^{72}$ from an on-device SMPL-X estimator as the
control variable in a closed loop with the array's precoder.  Two
regimes emerge from the same construction.  In the \emph{LOS
un-shadowing} regime, pose moves that take the user's body out of
the line of sight increase rate and decrease absorbed dose
simultaneously; we give a closed-form rate-gradient lemma and a
``win-win'' theorem.  In the \emph{body-as-passive-RIS} regime, pose
moves that orient the body to redirect specular paths trade rate for
dose along a body-universal Pareto slope $T_0/(1-T_0)\approx 0.32$
at \SI{28}{GHz}.  The two regimes are complementary, not
contradictory: LOS un-shadowing handles the open-window scenario,
the Pareto bound handles the multi-user crowd.  Differentiability of
the exposure operator $\Qabs(\thetab)$ in pose follows from
linear-blend-skinning smoothness and the Kirchhoff surface-integral
chain rule; the closed-loop bandwidth budget is
$\le$\,\SI{130}{kbps}, four orders of magnitude below the data rate
the loop negotiates over.  The hero claim is empirical and is left
to a companion experiment document.
\end{abstract}

\vspace{0.3em}
\noindent\textit{Style note.}~Like its siblings
\texttt{rib.tex} and \texttt{q\_complement.tex}, this is a wandering
companion, not a theorem-proof paper.  Theorems are theorems where
the proof is self-contained; propositions are claims I want a
reviewer to take or leave on sketches; conjectures are flagged.
The intended home for the deployable parts is \S\,IV of a JSAC2 SI
submission.

\tableofcontents
\bigskip

% ============================================================
\section{Why a fifth body-twin paper}\label{sec:why}
% ============================================================

The body-twin papers in the lineage \texttt{(paper, paper\_v2,
paper\_jsac\_v1\dots v5)} converged on a result whose body twin had
structurally shrunk to ``supply one scalar per slot.''  Under the
operational \mbox{5G NR FR2} per-stream MCS27 cap, regularised
zero-forcing on a tens-of-element panel saturates the rate cap
with $\sim$\SI{12}{dB} of SINR margin, and a one-line scaling
$\alpha=\sqrt{\eta_S\,\min_u L_{\mathrm{RL}}^{(u)}/\Pabs^{(u)}}$
enforces compliance.  Two-seed multi-seed cross-checks show the
elaborate dual-ascent solver and the closed-form scaling are
statistically tied on aggregate compliance.  The deployment value
of the twin over a worst-case Cauchy envelope is real but small.
The body twin is decorative.

The JSAC2 SI on Digital Twins for Wireless Networks asks for
something different.  Its scope keywords are \emph{closed-loop
optimization}, \emph{differentiable control},
\emph{twin-in-the-loop}, and \emph{cross-layer feedback}.  None of
these describe a one-line amplitude scaling.  We re-target the body
twin from constraint provider to control variable: the user's pose
is the channel-shaping degree of freedom the network negotiates
with, and the body twin is the model that makes the negotiation
tractable.

\subsection{Why now and why this user}

mmWave deployment slowed in part because (a) bodies and trees block
$25\text{--}40$\,dB at \SI{28}{GHz}, and (b) consumer-side
exposure concerns produce friction even where the dose is well
within ICNIRP limits.  Both fixes flow through the body twin.  The
network knows the user's body (\Cref{sec:diff} below); the user
sees a dose receipt for the rate they purchase (\Cref{sec:loop}
below).  The cost is one consent prompt per session; the upside is
both a capacity story and a trust story, with one piece of
machinery.

\subsection{What this gives up}

The v5 paper is a story the operator can run with no UE-side
cooperation at all.  Our spine assumes opt-in.  At opt-in fraction
$f$, the architecture's value to opted-out users is \emph{indirect}:
opted-in users' bodies still scatter coherently into the propagation
graph, and at plaza scale the cooperative-body-mediated path
contribution is order $f\cdot K\cdot\bar V_{\mathrm{cross}}$ where
$K$ is the local crowd size and $\bar V_{\mathrm{cross}}\sim 0.05$
is the average cross-body visibility (\texttt{rib.tex} \S\,11).  At
$f=5\%$, $K=50$, this is $\sim$0.13 cooperative paths per user, of
order\ a quarter-dB MIMO lift.  Below useful.  At $f=30\%$ the same
arithmetic gives $\sim$0.75 paths per user, an order of magnitude
larger and approaching the empirical rank-lift threshold of the
JSAC v5 plaza runs.  Scenario S4 of \texttt{scenarios.md} measures
the curve.

\subsection{The minimum technical content}

The construction needs three things to hold up to a careful
reviewer.  (i) The exposure operator must be differentiable in
$\thetab$ with a chain rule the device can compute
(\Cref{sec:diff}).  (ii) The capacity-vs-dose trade must be
characterised in both pose-move regimes — LOS un-shadowing
(\Cref{sec:los-unshadow}) and body-as-RIS (\Cref{sec:rib-pareto}) —
with the right Pareto bound in each.  (iii) The closed-loop must
have a bandwidth and latency budget compatible with a
phone-class compute envelope (\Cref{sec:loop}).  The rest is
framing.

% ============================================================
\section{Setup and notation}\label{sec:setup}
% ============================================================

\subsection{Geometry}

A single base station hosts an array of $M$ elements at known
positions $\{\rr_j\}_{j=1}^M$.  In the running examples we take a
$M=64$ uniform rectangular array (URA) with half-wavelength spacing
at carrier frequency $f_c=\SI{28}{GHz}$
($\lambda\approx\SI{1.07}{cm}$).  The array radiates the precoder
$\xvec\in\Complex^M$ with sum-power constraint
$\|\xvec\|^2\le P_{\mathrm{tx}}$.

A single user is described by a parametric body model.  We use
SMPL-X (Pavlakos et al.\ 2019) with shape vector
$\betab\in\Reals^{10}$ (identity, fixed once per user) and pose
vector $\thetab\in\Reals^{72}$ (root orientation plus 21 axis-angle
joint rotations plus a few facial / hand DOF that we ignore).  The
body mesh is the closed orientable surface
\begin{equation}\label{eq:mesh}
  \Sigma(\thetab,\betab)
   =\Mmesh\bigl(\LBS(\thetab,\betab,\mathbf{T})\bigr),
\end{equation}
with $\mathbf{T}\in\Reals^{V\times 3}$ the SMPL-X template
($V=\num{10475}$ vertices, $F=\num{20908}$ triangles), $\LBS$ the
linear blend skinning operator, and $\Mmesh$ the topology assembly.

The user holds a smartphone receiver at position $\rr_p(\thetab)
\in\Reals^3$.  The phone is anchored at the right wrist, which is
itself an SMPL-X joint, so $\rr_p$ inherits LBS smoothness in
$\thetab$.  The phone-to-body distance $d_p\equiv\|\rr_p-\rr_b\|$
where $\rr_b$ is the body centroid is $\sim$\SI{30}{cm} in the
running examples, deep in the body's near field at \SI{28}{GHz}.
This forces a near-field correction in \Cref{sec:nearfield}.

\subsection{Path-space factorisation reminder}

The propagation between the array and any spatial probe is
described by $N$ multipath contributions
$\{j(n),\khat_n,\bm\psi_n,\tau_n\}_{n=1}^N$ as in the JSAC paper:
$j(n)\in\{1,\dots,M\}$ the source element, $\khat_n$ the unit
propagation direction at the probe, $\bm\psi_n\in\Complex^3$ the
polarisation amplitude, and $\tau_n$ the path delay.  At the body,
the per-element steering response across the lit area factors as
$\Jmat^T\mathbf{a}(\khat_n)$, where $\Jmat\in\Complex^{M\times M}$
encodes element-to-element phase relationships and
$\mathbf{a}(\khat)$ is the array manifold.  The two operators that
recur:
\begin{align}
  \Qabs(\thetab)
  &= \int_{\Sigma(\thetab)}
     \Gtil(\rr;\thetab)^H\,\Gtil(\rr;\thetab)\,
     \mu(\rr;\thetab)\,\diff A(\rr),
  \label{eq:Qabs-def}\\[0.2em]
  \Qbist(\thetab,\kout)
  &= \frac{1}{Z_0}\,\Frib(\thetab,\kout)^H\,
     \Frib(\thetab,\kout),
  \label{eq:Qbist-def}
\end{align}
with $\mu(\rr;\thetab)=\nhat(\rr;\thetab)\cdot(-\khat)$ the
incidence cosine, $Z_0$ the free-space impedance, and
$\Frib(\thetab,\kout)$ the body scattering matrix from
\texttt{rib.tex} \S\,4 (we drop the $\betab$-dependence everywhere
since shape is fixed once per user).

For one user and one precoder the absorbed-power scalar and the
body-mediated contribution to receive power are
\begin{equation}\label{eq:Pabs-Pbody}
  \Pabs(\thetab,\xvec)=\xvec^H\Qabs(\thetab)\xvec,
  \qquad
  \Pbody(\thetab,\xvec)
   =\frac{1}{(4\pi d_{u\to\mathrm{UE}})^2}\,
    \xvec^H\Qbist(\thetab,\kout_{u\to\mathrm{UE}})\xvec,
\end{equation}
where $\kout_{u\to\mathrm{UE}}$ is the body-to-phone direction.  We
will use $\Pabs$ and $\Pbody$ throughout to keep notation uniform.

\subsection{What this note adds to its siblings}

Both operators were treated as static functions of pose in the JSAC
v1--v5 papers: pose was an estimated parameter, the twin
re-evaluated $\Qabs$ at the new estimate every \SI{100}{ms}.  In
this note, $\thetab$ is a decision variable.  The mathematics is
the same; the question posed of it is different.

% ============================================================
\section{The user as a reconfigurable surface}\label{sec:user}
% ============================================================

\subsection{The body-state vector and its actuation manifold}

The user control vector is
$\bm\varphi=(\thetab,\rr_{\mathrm{root}},\bm\omega_{\mathrm{root}})
\in\Reals^{78}$ (SMPL-X pose plus root translation plus root
orientation).  Not all $78$ DOF are negotiable at the same cost.
We adopt the four-tier classification of biomechanics ergonomic
scoring (RULA, OWAS, REBA): \emph{free} (eyebrow, finger, jaw —
no metabolic cost, no mmWave effect), \emph{comfortable} (head
yaw/pitch, torso yaw, arm and forearm joint angles),
\emph{costly} (spine bend, kneeling, leaning beyond support
polygon), \emph{rigid} (root translation in world coordinates).

We model the cost of the pose move from baseline $\thetab^{(0)}$
to proposed $\thetab$ as a separable convex
\begin{equation}\label{eq:movement-cost}
  C(\thetab-\thetab^{(0)})
   = \sum_{i=1}^{72}c_i\,(\theta_i-\theta_i^{(0)})^2
   + \sum_{i=1}^{72}\beta_i\,
       \log\!\bigl(1-(\theta_i/\theta_i^{\max})^2\bigr)^{-1},
\end{equation}
with $c_i$ small for comfortable DOF, large for costly DOF, and
$\beta_i$ the joint-limit log-barrier coefficient.  In the
companion experiment we calibrate $c_i$ such that $C\le 1$
corresponds to a \SI{5}{\degree} torso rotation, equivalent to a
RULA score increment of $1$ (\emph{low load}).  A \SI{30}{\degree}
turn is $C\approx 36$, RULA $3$ (\emph{significant load}); these
are the moves the gradient should not suggest.

\subsection{Update cadence}

Internal pose drift (breathing, postural sway) gives
$\mathrm{std}(\thetab)\sim$ a few tenths of a degree per joint at
the torso across \SI{100}{ms} windows.  This is the cadence the
\texttt{virtual\_imu.py} pipeline tracks at; the JSAC v5 machinery
budgets the twin against it unchanged.

Negotiated pose moves (the BS suggests, the user takes) complete in
$\SI{200}{ms}$ to $\SI{2}{s}$.  The control loop runs at
\SI{1}{Hz}.  The inner-precoder cadence is the per-slot
\SI{0.1}{ms}; pose-gradient evaluation runs at a few \SI{}{Hz} on a
phone-class GPU (cost analysis in \Cref{sec:diff}).

% ============================================================
\section{Differentiability of the body channel}\label{sec:diff}
% ============================================================

\subsection{LBS is smooth in pose, almost everywhere}

\begin{lemma}[LBS smoothness]\label{lem:lbs-smooth}
For fixed shape $\betab$, the map $\thetab\mapsto\Sigma(\thetab,
\betab)$ in \eqref{eq:mesh} is piecewise-smooth on $\Reals^{72}$.
Each vertex position $\rr_v(\thetab)$ is given by
\begin{equation}\label{eq:lbs-eq}
  \rr_v(\thetab)
  =\sum_{j=1}^{J}w_{vj}\bigl(\mathbf{R}_j(\thetab)\,
                              (\mathbf{T}_v+\mathbf{B}_v(\betab)
                                              -\mathbf{j}_j)
                              +\mathbf{j}_j+\mathbf{t}_j(\thetab)\bigr),
\end{equation}
where $\mathbf{R}_j(\thetab)\in SO(3)$ is the world rotation of
joint $j$, polynomial in $\sin\theta_{j,\cdot},\cos\theta_{j,\cdot}$;
$\mathbf{j}_j$ is the canonical joint position; and
$\mathbf{t}_j(\thetab)$ is the joint's forward-kinematic
translation.  Both are $C^\infty$ in $\thetab$.  Triangle normals
$\nhat_t=\mathrm{normalise}((\rr_2-\rr_1)\times(\rr_3-\rr_1))$ are
$C^\infty$ wherever the triangle is non-degenerate.
\end{lemma}

\subsection{Smoothing the silhouette}\label{sec:gelu}

The Heaviside visibility gate $H(\mu)$ in \eqref{eq:Qabs-def} is the
only non-smooth ingredient.  We replace it by the GELU smoothing of
TAP \S\,2.6:
\begin{equation}\label{eq:gelu}
  H_\epsilon(\mu)=\half\bigl(1+\mathrm{erf}(\mu/(\epsilon\sqrt{2}))\bigr),
\end{equation}
with $\epsilon=0.017$ chosen so the smoothing zone width is one
incidence-degree.  At Thelonious mesh size and \SI{28}{GHz} the
bias on $\tr\Qabs$ is $\le$\,$0.4\%$, well below the JSAC paper's
$\sim$5\% calibration floor.

\subsection{The end-to-end Jacobian}

\begin{theorem}[Pose differentiability of $\Qabs$]\label{thm:Q-diff}
For fixed shape and path set, with GELU-smoothed visibility, the
map $\thetab\mapsto\Qabs(\thetab)\in\Complex^{M\times M}$ is
$C^\infty$ on the LBS smooth set.  The Jacobian factorises as
\begin{empheq}[box=\tcbhighmath]{equation}\label{eq:Q-jac}
  \frac{\partial\Qabs}{\partial\theta_i}
  \;=\;\sum_{t=1}^{F}\sum_{n=1}^{N}\sum_{v\in t}
        \frac{\partial q_{t,n}}{\partial\rr_v}\cdot
        \frac{\partial\rr_v}{\partial\theta_i},
\end{empheq}
where $q_{t,n}$ is the per-(triangle, path) integrand (a smooth
function of vertex positions and path direction) and the LBS
Jacobian $\partial\rr_v/\partial\theta_i$ is exactly the SMPL-X
PyTorch autograd output.
\end{theorem}

\noindent Compose \cref{lem:lbs-smooth} with the smoothness of the
integrand (Fresnel coefficient is rational in $\mu$, area is
polynomial in vertices, GELU is $C^\infty$).  Sum the
$F\cdot N\cdot 3$ chain-rule contributions.  The same theorem
holds for $\Qbist(\thetab,\kout)$ verbatim with the additional
phase factor $e^{-ik_0\kout\cdot\rr}$.  $\square$

\subsection{The dB-per-degree LOS rate gradient}\label{sec:rate-grad}

The empirical hero claim depends on how steeply the rate changes
under small pose moves.  In the body-shadowed LOS regime, this is
dominated by un-shadowing rather than by changes in the bistatic
operators.  We give the closed form.

Let $v(\thetab)\in[0,1]$ be the LOS path's amplitude transmission
through the body.  $v=1$ when the LOS bypasses the body entirely;
$v\approx 0.06$ when the body fully blocks (\SI{25}{dB} body
loss).  Let
$h_{\mathrm{LOS}}(\thetab)=v(\thetab)\,h_{\mathrm{LOS},0}\,e^{i\phi_0}$
where $h_{\mathrm{LOS},0}$ is the unblocked LOS amplitude.

\begin{lemma}[Rate gradient in body-shadowed LOS]\label{lem:rate-grad}
Below the MCS cap, with single-stream MRT-MMSE precoding,
\begin{equation}\label{eq:rate-grad-los}
  \frac{\partial R}{\partial\theta_i}
  =\frac{B}{\ln 2}\,
    \frac{2\,\SINR_{\max}\,v(\thetab)\,(\partial v/\partial\theta_i)}
         {1+\SINR_{\max}\,v(\thetab)^2},
\end{equation}
where $\SINR_{\max}=P_{\mathrm{tx}}|h_{\mathrm{LOS},0}|^2/N_0B$ is
the unblocked-LOS SINR.  In the binding regime
$\SINR_{\max}\,v(\thetab)^2\ll 1$, this simplifies to
\begin{equation}\label{eq:rate-grad-binding}
  \frac{\partial R}{\partial\theta_i}\approx
  \frac{2B\,\SINR_{\max}\,v(\thetab)}{\ln 2}\,
        \frac{\partial v}{\partial\theta_i},
\end{equation}
which is linear in $v$ and in the un-shadowing rate
$\partial v/\partial\theta_i$.
\end{lemma}

\noindent Differentiate
$R=B\log_2(1+\SINR_{\max}v(\thetab)^2)$.  The rate gradient is
governed entirely by $\partial v/\partial\theta_i$, the
geometric-blockage derivative.

\begin{remark}[Magnitude estimate]\label{rem:dB-per-deg}
For a torso of width $w_b\approx\SI{30}{cm}$ at distance $r_b
\approx\SI{30}{m}$ from the BS, with the spine offset from the
body's lateral center by $\delta_s\approx\SI{8}{cm}$, a torso
rotation of $\delta\theta=\SI{1}{\degree}$ moves the chest's
lateral position relative to the LOS by $\delta_s\cdot\sin
(\delta\theta)\approx\SI{1.4}{mm}$.  This is $\sim$$\lambda/8$;
within the first Fresnel zone ($\sqrt{\lambda r_b}\approx
\SI{60}{cm}$) it modulates $v(\thetab)$ smoothly with a slope
$\partial v/\partial\theta\sim\delta_s/w_b\approx 0.27$ per radian,
$\sim 0.005$ per degree.  Combined with the binding-regime
expression \eqref{eq:rate-grad-binding} and $\SINR_{\max}\sim
\SI{30}{dB}$ at unblocked LOS, the per-degree rate gradient
peaks at $\sim$\SI{2}{dB/degree} of torso yaw at mid-shadow
($v\approx 0.5$).  For combined moves in three DOF the linearised
gain over a $C\le 1$ ball is therefore in the \SI{6}{dB} to
\SI{12}{dB} range.  This is the pre-experiment estimate;
\texttt{scenarios.md} S1 measures it.
\end{remark}

\subsection{Cost of one Jacobian-vector product}

The chain rule in \eqref{eq:Q-jac} costs one Fresnel-side
evaluation per (triangle, path) and one LBS-side evaluation per
(vertex, pose DOF).  In numbers: $F\cdot N\sim 10^6$ Fresnel-side
terms (each a $3\times M$ complex matrix, $\sim M^2\cdot 10^6=
4\cdot 10^9$ flops at $M=64$); $V\cdot 72\sim\num{750000}$
LBS-side terms (each a $3\times 1$ real vector, trivial).  The
forward pass is identical in cost to the JSAC paper's $\Qabs$
build, $\sim$\SI{7.5}{ms} per body on a single GPU at SMPL-X mesh
size.  A reverse-mode VJP of one scalar against $\thetab\in
\Reals^{72}$ costs $\sim$\SI{15}{ms} on a phone-class GPU.

% ============================================================
\section{Two pose-move regimes}\label{sec:regimes}
% ============================================================

The empirical and theoretical content of the spine splits cleanly
in two.  The pose moves available to the user fall in one of two
regimes, governed by which channel component dominates the rate
gradient.  We treat each in its own section.

\subsection{The dichotomy}

Define the LOS un-shadowing fraction
$v(\thetab)\in[0,1]$ and the body-mediated coupling
$\Pbody(\thetab,\xvec)$ as in \eqref{eq:Pabs-Pbody}.  The total
receive power is
\begin{equation}\label{eq:Prx}
  P_{\mathrm{rx}}(\thetab,\xvec)
  =P_{\mathrm{LOS}}(\thetab,\xvec)+\Pbody(\thetab,\xvec)
  +\text{wall-multipath terms}.
\end{equation}
A pose move $\thetab\to\thetab+\delta\thetab$ changes both
contributions.  Define the dominant component:
\begin{itemize}[leftmargin=1.6em,itemsep=2pt]
\item \textbf{Regime~A (LOS un-shadowing).}  $|\partial
P_{\mathrm{LOS}}/\partial\thetab|\gg|\partial\Pbody/\partial
\thetab|$.  The pose move is moving the body in or out of the LOS
path.
\item \textbf{Regime~B (body as RIS).}  $|\partial\Pbody/\partial
\thetab|\gtrsim|\partial P_{\mathrm{LOS}}/\partial\thetab|$.  The
pose move is re-orienting the body to redirect specular paths.
\end{itemize}
The two regimes need different theoretical instruments.

% ============================================================
\section{Regime A: LOS un-shadowing}\label{sec:los-unshadow}
% ============================================================

In Regime~A the pose move's rate effect is captured by
\Cref{lem:rate-grad}.  The dose effect runs in the opposite
direction.

\subsection{Dose decreases as LOS un-shadows}

When the body un-shadows the LOS, the lit-area integral in
\eqref{eq:Qabs-def} shrinks.  The geometry is intuitive: the same
projected area that intercepts the LOS is the area that absorbs the
LOS power.  Concretely, for a single LOS path with amplitude
$|h_{\mathrm{LOS},0}|^2$ and visibility $v(\thetab)$,
\begin{equation}\label{eq:Pabs-LOS}
  \Pabs^{\mathrm{LOS}}(\thetab,\xvec)
  =(1-v(\thetab)^2)\,|h_{\mathrm{LOS},0}|^2\,P_{\mathrm{tx}},
\end{equation}
since the absorbed fraction is the complement of the transmitted
fraction.  Thus
$\partial\Pabs^{\mathrm{LOS}}/\partial\theta_i = -2v(\thetab)\,
(\partial v/\partial\theta_i)\,|h_{\mathrm{LOS},0}|^2P_{\mathrm{tx}}$.

\subsection{The win-win theorem}

\begin{theorem}[LOS un-shadowing is Pareto-dominant]\label{thm:los-pareto}
In Regime~A, every pose move with
$\partial v/\partial\theta_i>0$ simultaneously
\emph{(i)}~increases the rate
$\partial R/\partial\theta_i>0$
(\Cref{lem:rate-grad}), and
\emph{(ii)}~decreases the absorbed dose
$\partial\Pabs/\partial\theta_i<0$
\eqref{eq:Pabs-LOS}.  Below the MCS cap, the trade is monotone in
$v$: rate is concave-increasing in $v$, dose is concave-decreasing
in $v$.
\end{theorem}

\noindent Both signs follow directly from \eqref{eq:rate-grad-los}
and \eqref{eq:Pabs-LOS}.  The concavity statements follow from
$R\sim\log(1+v^2)$ (concave in $v^2$) and $\Pabs^{\mathrm{LOS}}\sim
(1-v^2)$ (linear in $v^2$, hence concave-decreasing in $v$).
$\square$

\Cref{thm:los-pareto} is the theoretical backbone of the hero
scenario.  In Regime~A there is no Pareto trade: pose moves that
unshadow the LOS are strictly better on both axes.  The dose
receipt the user sees is honest about this: ``you bought
\SI{12}{dB} of rate and reduced your dose by \SI{8}{dB} in the
process.''  This is the friendliest possible UX framing of an
exposure-aware loop.

\subsection{The MCS-cap saturation plateau}

Above the MCS cap, the rate gradient vanishes:
\begin{equation}\label{eq:cap-zero}
  R(\thetab)=R_{\max}\;\;\Rightarrow\;\;
  \frac{\partial R}{\partial\theta_i}=0,
\end{equation}
so the loss \eqref{eq:loss} below has no rate term to maximise.
The dose term remains active, however, and the gradient on
$\thetab$ now points only at minimising $\Pabs$.

\begin{corollary}[Operating point at saturation]\label{cor:cap}
At a pose where the SINR has reached the MCS cap, the local optimum
of \eqref{eq:loss} is the lowest-dose pose in the comfort
neighbourhood that still attains the cap.  Equivalently, the
gradient ascends the iso-rate set's lowest-dose direction, which
in Regime~A is the same direction as further LOS un-shadowing.
\end{corollary}

The Pareto plot in this regime has a vertical edge at $R=R_{\max}$:
all the gain along that edge is dose reduction at constant rate.
This is the engineering statement the BS sends to the UE: ``you
have already maxed out your rate; further pose moves only reduce
your dose receipt.''

% ============================================================
\section{Regime B: body as passive RIS}\label{sec:rib-pareto}
% ============================================================

In Regime~B the LOS is either unavailable (NLOS) or already
saturated (above-cap).  Pose moves no longer trade in
LOS un-shadowing; they reorient the body so that body-mediated
specular paths reach the receiver.  The Pareto bound of
\texttt{rib.tex} \S\,8 governs this regime, with the
single-user pose-conditional refinement below.

\subsection{The single-user Pareto bound, fixed}

\begin{theorem}[Single-user pose-conditional Pareto, trace form]
\label{thm:pareto-1}
Under the assumptions of \texttt{rib.tex} \S\,3 (high-frequency
PO, far-field observer, single bounce, pseudo-Brewster), at any
fixed pose $\thetab$, at receiver direction
$\kout=\kout_{u\to\mathrm{UE}}(\thetab)$, and at any precoder
$\xvec$,
\begin{empheq}[box=\tcbhighmath]{equation}\label{eq:pareto-1}
  \xvec^H\Qbist(\thetab,\kout)\xvec
  \;\le\;
  \frac{1-T_0}{T_0}\,\xvec^H\Qabs^{(\kout)}(\thetab)\xvec
  \;+\;O\!\bigl(5\%\bigr),
\end{empheq}
with $\Qabs^{(\kout)}(\thetab)$ the $\kout$-visible-area-restricted
absorbed-power operator.  Equality holds when $\xvec$ aligns with
the top eigenvector of $\Qabs^{(\kout)}$.  Taking traces and
bounding by visibility,
\begin{equation}\label{eq:pareto-1-trace}
  \tr\Qbist(\thetab,\kout)
  \;\le\;
  \frac{1-T_0}{T_0}\,\mathcal{V}_{\mathrm{out}}(\thetab)\,
                       \tr\Qabs(\thetab),
\end{equation}
where $\mathcal{V}_{\mathrm{out}}(\thetab)\in[0,1]$ is the trace
fraction $\tr\Qabs^{(\kout)}/\tr\Qabs$.  The body-mediated power
contribution to the receive flux is bounded by
\begin{equation}\label{eq:Pbody-bound}
  \Pbody(\thetab,\xvec)
  \;\le\;
  \frac{1}{(4\pi d_{u\to\mathrm{UE}})^2}\,
   \frac{1-T_0}{T_0}\,\mathcal{V}_{\mathrm{out}}(\thetab)\,
   M\,\Pabs(\thetab,\xvec)/\tr\Qabs,
\end{equation}
i.e.\ for any normalised precoder, $\Pbody\le c(\thetab)\Pabs$ with
the body-and-geometry constant
$c(\thetab)=\frac{1-T_0}{T_0}\mathcal{V}_{\mathrm{out}}(\thetab)/
(4\pi d_{u\to\mathrm{UE}})^2$.
\end{theorem}

\noindent The operator inequality \eqref{eq:pareto-1} is
\texttt{rib.tex} \S\,7 (per-direction pseudo-Brewster locking)
applied to the $\kout$-restricted operators on both sides.  The
trace inequality \eqref{eq:pareto-1-trace} is its trace; the
visibility factor $\mathcal{V}_{\mathrm{out}}\le 1$ accounts for the
restricted integration domain.  $\square$

\begin{remark}\label{rem:fixed-op-ineq}
The earlier draft erroneously asserted an operator inequality
$\Qabs^{(\kout)}\preceq\mathcal{V}_{\mathrm{out}}\Qabs$, which is
not implied by domain restriction (the eigenvalues of
$\Qabs^{(\kout)}$ can be at most equal to those of $\Qabs$ but are
not bounded by a trace-fraction prefactor).  The correct statement
is the trace bound \eqref{eq:pareto-1-trace}, which is what is
used in the corollaries below.
\end{remark}

\subsection{The gait-action principle, restated}

\begin{corollary}[Gait-action, body-mediated component]\label{cor:gait}
At any precoder $\xvec$ aligned with the top eigenvector of
$\Qabs^{(\kout)}$, the ratio of body-mediated capacity-per-dose,
$\Pbody/\Pabs$, is bounded by $c(\thetab)$ from
\eqref{eq:Pbody-bound} and is maximised when $\xvec$ aligns with
that eigenvector.  Pose moves that scale
$\mathcal{V}_{\mathrm{out}}(\thetab)$ scale this maximum
proportionally; the body-universal slope $(1-T_0)/T_0$ is
unchanged.

The total capacity-per-dose ratio is \emph{not} bounded by
$c(\thetab)$, because the LOS contribution
$P_{\mathrm{LOS}}/\Pabs$ has its own behaviour governed by
\Cref{thm:los-pareto}.
\end{corollary}

This restated form fixes the over-reach in v1.  The body-mediated
component is governed by the body-universal Pareto slope; the LOS
component is governed by the win-win theorem of Regime~A.

\subsection{Pose-noise-aware bound}\label{sec:pose-noise}

The pose estimator on the device has covariance $\bm\Sigma_\theta
\succeq 0$.  At evaluation time, the BS uses the mean estimate
$\hat\thetab$ rather than the true $\thetab$.  Linearising the
Pareto bound around $\hat\thetab$,

\begin{proposition}[Noise-aware Pareto]\label{prop:noisy-pareto}
For pose noise with covariance $\bm\Sigma_\theta$, the expected
body-mediated power is bounded by
\begin{equation}\label{eq:noisy-pareto}
  \mathbb{E}_{\thetab\sim\mathcal{N}(\hat\thetab,\bm\Sigma_\theta)}
   \bigl[\Pbody(\thetab,\xvec)\bigr]
  \;\le\;
  c(\hat\thetab)\,\Pabs(\hat\thetab,\xvec)
  +\half\,\tr\!\bigl(\bm\Sigma_\theta\,\nabla^2_\thetab
                     \,\bigl(c\Pabs\bigr)\big|_{\hat\thetab}\bigr)
  +O(\|\bm\Sigma_\theta\|^{3/2}).
\end{equation}
\end{proposition}

\noindent Taylor-expand $\Pbody$ around $\hat\thetab$, take
expectations.  The first-order term zeroes; the second-order term
is the Hessian-trace correction.  $\square$

\Cref{prop:noisy-pareto} says the noisy-pose bound differs from
the noiseless one by a Hessian-trace term that is
$O(\|\bm\Sigma_\theta\|)$.  At the JSAC v5 IMU's $\sim
\SI{0.5}{\degree}$ RMS, this term is $\sim 1\%$ of the leading
order.  The pose noise's effect on the Pareto bound is
sub-percent.

% ============================================================
\section{Near-field correction for body-to-phone}\label{sec:nearfield}
% ============================================================

The phone-to-body distance $d_p\approx\SI{30}{cm}$ is well inside
the body's Fraunhofer distance $D^2/\lambda\approx\SI{300}{m}$ at
\SI{28}{GHz}.  $\Frib(\thetab,\kout)$ as defined in
\texttt{rib.tex} \S\,4 is a far-field object; using it directly to
predict body-to-phone amplitude is wrong by the Fresnel-Kirchhoff
correction.

\subsection{The near-field body scattering matrix}

For a receiver at finite position $\rr_p$, the far-field
$1/(4\pi R)$ in \texttt{rib.tex} eq.\ for $\bvec_{u\to\mathrm{UE}}$
is replaced by the Fresnel-Kirchhoff kernel
$e^{ik_0|\rr-\rr_p|}/(4\pi|\rr-\rr_p|)$.  The body scattering
matrix becomes
\begin{equation}\label{eq:Fb-NF}
  \Frib^{(n),\mathrm{NF}}(\thetab,\rr_p)
  =\frac{ik_0}{4\pi}\,\Pi^{\mathrm{out}}(\rr_p)\!
   \int_{\Sigma(\thetab)}\!\!\mathbf{R}(\rr;\khat_n)\,
   \frac{e^{ik_0(|\rr-\rr_p|-\khat_n\cdot\rr)}}{|\rr-\rr_p|}\,
   H(\mu_n)\,H(\nhat\cdot\hat{\bm\eta})\,\diff A,
\end{equation}
where $\hat{\bm\eta}=(\rr_p-\rr)/|\rr_p-\rr|$ is the per-point
direction to the phone (no longer a constant $\kout$).

\subsection{The near-field Pareto correction}

\begin{proposition}[Near-field Pareto]\label{prop:nf-pareto}
Under the same pseudo-Brewster regime as
\Cref{thm:pareto-1}, with the receiver in the body's Fresnel zone
at distance $d_p$ from the body centroid,
\begin{equation}\label{eq:nf-pareto}
  \Pbody^{\mathrm{NF}}(\thetab,\xvec)
  \;\le\;
  \frac{1-T_0}{T_0}\,
  \frac{\mathcal{V}_{\mathrm{out}}^{\mathrm{NF}}(\thetab)}
       {(4\pi d_p)^2}\,
   M\,\Pabs(\thetab,\xvec)/\tr\Qabs\,
  \cdot g_{\mathrm{NF}}(d_p,D),
\end{equation}
with the near-field correction factor $g_{\mathrm{NF}}(d_p,D)
=1+O(D/d_p)$ for $d_p<D^2/\lambda$.
\end{proposition}

\noindent Repeat the proof of \Cref{thm:pareto-1} with the
Fresnel-Kirchhoff kernel; the per-point direction $\hat{\bm\eta}$
breaks the constant-$\kout$ pseudo-Brewster locking but the local
locking applies pointwise; integrating gives the Hessian-trace
correction.  Quantification of $g_{\mathrm{NF}}$ for a body at
$d_p\approx\SI{30}{cm}$ is a numerical exercise (one Kirchhoff
integral per pose); we estimate $g_{\mathrm{NF}}\le 2$ for the
worst-case torso-to-hand geometry.  $\square$

\subsection{What this means for the spine}

The near-field correction modifies the Pareto slope by at most a
factor of $2$.  The body-universal $(1-T_0)/T_0=0.59$ at
\SI{28}{GHz} becomes a body-and-geometry-specific
$\le 1.2$ in the worst case.  This widens the trade but does not
break it.  The user-facing dose receipt remains honest; only the
displayed conversion rate changes.

The near-field correction is most relevant in Regime~B, where the
body-mediated path goes from the BS off the body to the phone.  In
Regime~A the phone is hit directly by the LOS and the near field
is irrelevant.  Hero scenario S1 of \texttt{scenarios.md} is in
Regime~A and the near-field correction does not affect its
analysis; future S3 (multi-user) is in Regime~B and does need it.

% ============================================================
\section{The joint loss and closed loop}\label{sec:loss}
% ============================================================

\subsection{Composite objective}

For $\lambda_E,\lambda_M>0$,
\begin{empheq}[box=\tcbhighmath]{equation}\label{eq:loss}
  \mathcal{L}(\thetab,\xvec;\thetab^{(0)})
  = -R(\thetab,\xvec)
    \;+\;\lambda_E\,\bigl[\Pabs(\thetab,\xvec)-\Lrl\bigr]_+
    \;+\;\lambda_M\,C(\thetab-\thetab^{(0)}).
\end{empheq}
$R$ is the MCS-capped achievable rate; $\Pabs$ is from
\eqref{eq:Pabs-Pbody}; $C$ is the comfort cost
\eqref{eq:movement-cost}; $[\cdot]_+$ is GELU-smoothed at the same
$\epsilon$ as the visibility for differentiability.

\subsection{Inner-outer alternation and convergence}

The natural solver alternates a closed-form precoder optimum with a
gradient step on $\thetab$.  By Danskin's envelope theorem,
\begin{equation}\label{eq:envelope}
  \frac{\diff R}{\diff\thetab}\bigg|_{\xvec^*(\thetab)}
  = \frac{\partial R}{\partial\thetab},
\end{equation}
so the outer-loop gradient does not need to differentiate through
the inner QCQP.

\begin{proposition}[Monotone descent]\label{prop:descent}
Under \cref{thm:Q-diff}, the alternating scheme produces a
non-increasing sequence
$\mathcal{L}(\thetab^{(\ell)},\xvec^{(\ell)})$ with limit a
stationary point.
\end{proposition}

\subsection{Convergence under partial compliance}\label{sec:compliance}

The user does not always take the suggested move.  Model the
acceptance as a probability $p_a(\delta\thetab^*)$ that decreases
with the magnitude of the suggestion:
\begin{equation}\label{eq:p-accept}
  p_a(\delta\thetab^*)=
  \exp\!\bigl(-\gamma\,C(\delta\thetab^*)\bigr),
  \qquad\gamma>0.
\end{equation}
Under partial compliance, the realised pose at slot $\ell+1$ is
$\thetab^{(\ell+1)}=\thetab^{(\ell)}+a_\ell\delta\thetab^*$ with
$a_\ell\sim\mathrm{Bernoulli}(p_a(\delta\thetab^*))$.

\begin{proposition}[Biased convergence]\label{prop:bias}
The expected stationary point of the alternating scheme under
acceptance model \eqref{eq:p-accept} is biased away from the
unconstrained optimum by an amount
$O\bigl((1-\bar p_a)\rho_{\mathrm{comfort}}\bigr)$, where
$\bar p_a$ is the time-averaged acceptance and
$\rho_{\mathrm{comfort}}$ is the comfort budget.
\end{proposition}

\noindent The expected pose update has magnitude $\bar p_a
\|\delta\thetab^*\|$.  Substituting into the descent inequality
gives a stationarity gap of $(1-\bar p_a)$ times the noiseless
gap.  $\square$

In practice $\bar p_a\sim 0.3\text{--}0.7$ in user studies of
notification-driven micro-tasks; the architecture survives at
small $\bar p_a$ because the rate gradient is small there
(\Cref{lem:rate-grad} is dominated by un-shadowing geometry, not
by the magnitude of the move).

% ============================================================
\section{Closed-loop architecture}\label{sec:loop}
% ============================================================

\subsection{Block diagram}

\begin{center}
\begin{tikzpicture}[
  node distance=10mm and 14mm, font=\small,
  block/.style={draw, rounded corners, align=center, minimum height=8mm,
                minimum width=18mm, fill=gray!8},
  arrow/.style={-{Stealth[length=2mm]}, thick},
]
\node[block] (bs) {BS array \\ M=64};
\node[block, right=of bs] (rt) {Path \\ dictionary};
\node[block, right=of rt] (ue) {UE phone \\ (SMPL-X)};
\node[block, below=of ue] (user) {User \\ (pose actuator)};
\node[block, left=of user] (imu) {IMU + AR \\ pose track};
\node[block, left=of imu] (qabs) {$\Qabs(\hat\thetab)$ \\ + gradient};
\node[block, below=of bs] (prec) {Precoder \\ ZF + projection};
\draw[arrow] (bs)--(rt) node[midway, above]{$\xvec$};
\draw[arrow] (rt)--(ue) node[midway, above]{paths};
\draw[arrow] (ue)--(user) node[midway, right]{suggest $\delta\thetab^*$};
\draw[arrow] (user)--(imu) node[midway, above]{move};
\draw[arrow] (imu)--(qabs) node[midway, above]{$\hat\thetab$};
\draw[arrow] (qabs)--(prec) node[midway, left]{$\nabla\mathcal{L}$};
\draw[arrow] (prec)--(bs) node[midway, left]{};
\draw[arrow] (qabs.north) |- (ue.west) node[midway, above]{accept/reject};
\draw[arrow, dashed] (rt) -- (qabs.north east);
\end{tikzpicture}
\end{center}

\noindent The dashed arrow is the path-dictionary downlink; the
solid arrows are the per-slot data flow.

\subsection{Information channels and rates}

Each path is described by 8 floats (direction $\khat\in S^2$:
2~floats; complex polarisation: 6 floats normalised; delay: 1
float; amplitude: 1 float).  At 32-bit precision and $N=50$ paths
re-evaluated at \SI{10}{Hz}, the path-dictionary downlink
amounts to $50\cdot 10\cdot 32\cdot 10=\SI{128}{kbps}$.

\begin{center}
\begin{tabular}{lll}
\toprule
\textbf{Channel} & \textbf{Payload} & \textbf{Rate} \\
\midrule
DL: path dictionary               & $50$ paths $\times$ 8 floats $\times$ \SI{10}{Hz}        & \SI{128}{kbps} \\
DL: pose suggestion (network-side) & $72$ floats $\times$ \SI{1}{Hz}                          & \SI{2.3}{kbps} \\
UL: pose estimate $\hat\thetab$    & $72$ floats $\times$ \SI{1}{Hz}                          & \SI{2.3}{kbps} \\
UL: accept/reject                 & $1$ bit $\times$ \SI{1}{Hz}                              & \SI{1}{bps} \\
UL: shape $\betab$ (one-time)     & $10$ floats once per session                            & \SI{40}{B} \\
\bottomrule
\end{tabular}
\end{center}

\noindent Total bidirectional control bandwidth is
$\le$\,\SI{135}{kbps}, four orders of magnitude below the data
rate the loop negotiates over.

\subsection{Privacy by gradient passing, with a caveat}\label{sec:privacy}

In the user-side variant, the phone transmits only:
\emph{(i)}~the current pose estimate, \emph{(ii)}~the suggestion
acceptance bit, and \emph{(iii)}~optionally an aggregated
absorbed-dose receipt.  Body shape $\betab$ stays on-device.  This
is structurally similar to federated learning: heavy data stays
local, summary statistics flow.

\begin{remark}[Acceptance-pattern inference]\label{rem:accept-infer}
Even with shape and pose held local, the BS sees the time-series
of acceptance bits and pose-estimate updates.  From this it can
infer the user's pose distribution, response latency, and likely
activity.  The strong privacy property is ``the operator does not
see body geometry,'' not ``the operator learns nothing about the
user.''  Differential-privacy treatments of the acceptance bit are
the obvious mitigation but are out of scope.
\end{remark}

\subsection{No new hardware}

Both endpoints run on commodity gear: the BS is a vanilla 5G FR2
panel with the JSAC v5 ECBF / projected-ZF precoder; the UE is a
smartphone with IMU and a pose-estimator.  The added compute on the
device is one VJP per second on the SMPL-X mesh (\SI{15}{ms} on
phone-class GPU per \Cref{thm:Q-diff}).  This is the same scale as
real-time face filters; the engineering is plumbing.

% ============================================================
\section{Connection to the JSAC2 SI scope}\label{sec:si}
% ============================================================

The SI keywords map to this construction in the following one-to-one
way.

\begin{enumerate}[leftmargin=1.6em,itemsep=2pt]
\item \emph{Closed-loop optimization.}  \Cref{sec:loop} is a literal
closed loop with measurable feedback and bounded latency.
\item \emph{Differentiable control.}  \Cref{thm:Q-diff} +
\Cref{lem:rate-grad}.
\item \emph{Twin-in-the-loop architectures.}  The body twin is
load-bearing; remove it and the gradient is ill-defined.
\item \emph{Cross-layer optimization.}  PHY (rate, exposure) couples
to biomechanical (pose, comfort).
\item \emph{Application-aware design.}  The user's session is the
application; the construction parameterises in it.
\item \emph{Predictive DT analytics.}  The phone displays the
current dose receipt and the predicted dose-and-rate trajectory.
\item \emph{Trustworthy data collection.}  User-side variant
(\Cref{sec:privacy}); body data stays local.
\item \emph{XR-adjacent.}  AR / VR sessions already track body
pose; the construction reuses that pipeline.  This is the weakest
of the eight; see also \Cref{sec:notthis}.
\end{enumerate}

% ============================================================
\section{What this is not}\label{sec:notthis}
% ============================================================

\begin{enumerate}[leftmargin=1.6em,itemsep=3pt]
\item Not a JCAS / ISAC paper.  We use on-device pose tracking, not
backscatter pose estimation.
\item Not a beamforming paper.  The precoder is the JSAC v5
projected-ZF closed form.
\item Not a pose-estimation paper.  The pose estimator is
borrowed from Madgwick / Mahony / IMU-SMPL-X literature.
\item Not an RIS paper.  The body is not actively tunable; gait is.
\item Not a regulatory-compliance paper.  Exposure constraints are
the engineering substrate; we do not certify a deployment.
\item Not a hardware paper.  No new hardware on either side.
\item Not an XR paper.  The XR connection is real but is one
sentence, not the spine.
\end{enumerate}

% ============================================================
\section{Loose threads and conjectures}\label{sec:threads}
% ============================================================

\begin{enumerate}[leftmargin=1.6em,itemsep=3pt]
\item \emph{Multi-user RIHB.}  Joint loss
$\mathcal{L}(\thetab_1,\thetab_2,\xvec)$.  Conjecture: a
Stackelberg-style lift exists when both users follow gradients.
Scenario S3 measures.
\item \emph{Adversarial poses.}  Worst-case pose for a given
precoder is the global maximum of $\Pabs(\thetab)$ over a
comfort-bounded ball.  Robustness question for the BS: how much
rate must it leave for safety?
\item \emph{Federated $\betab$.}  Population-level capacity
statistics without seeing individual $\betab$.
\item \emph{Sub-\SI{6}{GHz}.}  Pseudo-Brewster collapse holds at
mmWave; below \SI{6}{GHz} the body's $T_0$ is
polarisation-dependent and the slope is no longer body-universal.
\item \emph{Pose suggestion as channel coding.}  The BS suggests
a \emph{distribution} over poses with priors from the comfort
metric; rate-distortion view of human-in-the-loop control.
\item \emph{Comfort metric calibration.}  The weights $c_i$ in
\eqref{eq:movement-cost} need a small ergonomics user study to
calibrate.  RULA / REBA give the framework; specific numbers do
not exist for our pose vocabulary.
\item \emph{Negotiation as a game.}  Operator that bills on data
delivered while user pays in pose deviation: two-player game with
known utilities.  Conjecturally the loop converges to the Nash
equilibrium.
\item \emph{Quantitative near-field.}  $g_{\mathrm{NF}}(d_p,D)$ in
\Cref{prop:nf-pareto} is bounded but not computed.  One Kirchhoff
integral per pose suffices.
\item \emph{The role of clothing.}  Textile roughness perturbs the
specular peak by $\sim 30\%$ at \SI{28}{GHz}; the
JSAC v5 calibration pipeline absorbs this.  Whether the absorbed
fraction is consistent across pose (i.e.\ the calibration is
pose-stationary) is open and is a real engineering question.
\item \emph{Foundation-model dose receipt.}  The dose receipt the
user sees could be augmented by a per-tissue health-impact
predictor.  Is well outside scope; flagged as a future-work
sentence.
\end{enumerate}

% ============================================================
\section{Summary}\label{sec:sum}
% ============================================================

The body twin in the JSAC paper is a noun; this note makes it a
verb.  The same surface-integral operators that already estimate
dose now flow gradients to the user, who is free to take or ignore
them.  Two regimes emerge cleanly: in LOS un-shadowing
(\Cref{thm:los-pareto}) pose moves are Pareto-dominant, increasing
rate and decreasing dose simultaneously; in body-as-RIS
(\Cref{thm:pareto-1}) pose moves trade rate for dose along a
body-universal slope $T_0/(1-T_0)\approx 0.32$ at \SI{28}{GHz}.
The differentiability contract (\Cref{thm:Q-diff}) is small; the
information budget of the loop ($\le$\,\SI{135}{kbps}) is
negligible; the convergence under partial compliance
(\Cref{prop:bias}) is graceful.  What remains is empirical: how
much capacity does a few-degree pose move buy in the body-shadowed
LOS regime, and at what dose.  The companion experiment in
\texttt{scenarios.md} S1 answers that.

\end{document}
